\documentclass[10pt,a4paper]{amsart}
\usepackage{amsmath,amssymb,amsthm,graphicx}
% Retained minimal preamble for the Stacks Project chapter sites-cohomology.tex (Cohomology on Sites).
% Verbatim excerpt of the chapter preamble: Xy-pic setup (native preamble lines 14-18) and the native
% theorem declarations with the mathematical macros (native preamble lines 154-217). Copied mechanically;
% nothing added, removed or redefined. The wrapper class and the amsmath/amssymb/amsthm/graphicx base
% packages are supplied by the assembler, so the class/externaldocument part of the native preamble is omitted.
% For commutative diagrams we use Xy-pic
\usepackage[all]{xy}

% We use 2cell for 2-commutative diagrams.
\xyoption{2cell}
% Theorem environments.
%
\theoremstyle{plain}
\newtheorem{theorem}[subsection]{Theorem}
\newtheorem{proposition}[subsection]{Proposition}
\newtheorem{lemma}[subsection]{Lemma}

\theoremstyle{definition}
\newtheorem{definition}[subsection]{Definition}
\newtheorem{example}[subsection]{Example}
\newtheorem{exercise}[subsection]{Exercise}
\newtheorem{situation}[subsection]{Situation}

\theoremstyle{remark}
\newtheorem{remark}[subsection]{Remark}
\newtheorem{remarks}[subsection]{Remarks}

\numberwithin{equation}{subsection}

% Macros
%
\def\lim{\mathop{\mathrm{lim}}\nolimits}
\def\colim{\mathop{\mathrm{colim}}\nolimits}
\def\Spec{\mathop{\mathrm{Spec}}}
\def\Hom{\mathop{\mathrm{Hom}}\nolimits}
\def\Ext{\mathop{\mathrm{Ext}}\nolimits}
\def\SheafHom{\mathop{\mathcal{H}\!\mathit{om}}\nolimits}
\def\SheafExt{\mathop{\mathcal{E}\!\mathit{xt}}\nolimits}
\def\Sch{\mathit{Sch}}
\def\Mor{\mathop{\mathrm{Mor}}\nolimits}
\def\Ob{\mathop{\mathrm{Ob}}\nolimits}
\def\Sh{\mathop{\mathit{Sh}}\nolimits}
\def\NL{\mathop{N\!L}\nolimits}
\def\CH{\mathop{\mathrm{CH}}\nolimits}
\def\proetale{{pro\text{-}\acute{e}tale}}
\def\etale{{\acute{e}tale}}
\def\QCoh{\mathit{QCoh}}
\def\Ker{\mathop{\mathrm{Ker}}}
\def\Im{\mathop{\mathrm{Im}}}
\def\Coker{\mathop{\mathrm{Coker}}}
\def\Coim{\mathop{\mathrm{Coim}}}

% Boxtimes
%
\DeclareMathSymbol{\boxtimes}{\mathbin}{AMSa}{"02}

%
% Macros for moduli stacks/spaces
%
\def\QCohstack{\mathcal{QC}\!\mathit{oh}}
\def\Cohstack{\mathcal{C}\!\mathit{oh}}
\def\Spacesstack{\mathcal{S}\!\mathit{paces}}
\def\Quotfunctor{\mathrm{Quot}}
\def\Hilbfunctor{\mathrm{Hilb}}
\def\Curvesstack{\mathcal{C}\!\mathit{urves}}
\def\Polarizedstack{\mathcal{P}\!\mathit{olarized}}
\def\Complexesstack{\mathcal{C}\!\mathit{omplexes}}
% \Pic is the operator that assigns to X its picard group, usage \Pic(X)
% \Picardstack_{X/B} denotes the Picard stack of X over B
% \Picardfunctor_{X/B} denotes the Picard functor of X over B
\def\Pic{\mathop{\mathrm{Pic}}\nolimits}
\def\Picardstack{\mathcal{P}\!\mathit{ic}}
\def\Picardfunctor{\mathrm{Pic}}
\def\Deformationcategory{\mathcal{D}\!\mathit{ef}}

% --- appended 2026-10-12 (Chief): the project own definitions of the reference/slogan/history
% display environments, copied mechanically from the shared Stacks preamble lines 8-12
% (\usepackage{verbatim} plus the three \newenvironment lines, which the chapter preamble keeps
% verbatim from the project preamble). Nothing added, removed or redefined: these environments
% carry no mathematical content and the project itself typesets them as comments. Without these
% three lines a counted statement whose own text uses them fails with Environment slogan undefined.
\usepackage{verbatim}
\newenvironment{reference}{\comment}{\endcomment}
\newenvironment{slogan}{\comment}{\endcomment}
\newenvironment{history}{\comment}{\endcomment}
\usepackage{cleveref}
\crefname{originallemma}{original lemma}{original lemmas}
\Crefname{originallemma}{Original lemma}{Original lemmas}
\makeatletter
\newlabel{topology-lemma-closed-in-quasi-compact}{{005C}{source}{}{sourcepointer-1}{}}
\newlabel{topology-lemma-closed-in-quasi-compact@cref}{{[originallemma][0][]005C}{[1][0][]source}{}{}{}}
\makeatother
\crefname{originallemma}{original lemma}{original lemmas}
\Crefname{originallemma}{Original lemma}{Original lemmas}
\makeatletter
\newlabel{topology-lemma-fibre-product-closed}{{08ZH}{source}{}{sourcepointer-2}{}}
\newlabel{topology-lemma-fibre-product-closed@cref}{{[originallemma][0][]08ZH}{[1][0][]source}{}{}{}}
\makeatother
\crefname{originallemma}{original lemma}{original lemmas}
\Crefname{originallemma}{Original lemma}{Original lemmas}
\makeatletter
\newlabel{topology-lemma-locally-quasi-compact-Hausdorff}{{08ZR}{source}{}{sourcepointer-3}{}}
\newlabel{topology-lemma-locally-quasi-compact-Hausdorff@cref}{{[originallemma][0][]08ZR}{[1][0][]source}{}{}{}}
\makeatother
\crefname{originalsection}{original section}{original sections}
\Crefname{originalsection}{Original section}{Original sections}
\makeatletter
\newlabel{topology-section-Hausdorff}{{08ZD}{source}{}{sourcepointer-4}{}}
\newlabel{topology-section-Hausdorff@cref}{{[originalsection][0][]08ZD}{[1][0][]source}{}{}{}}
\makeatother
\crefname{originaltheorem}{original theorem}{original theorems}
\Crefname{originaltheorem}{Original theorem}{Original theorems}
\makeatletter
\newlabel{topology-theorem-tychonov}{{08ZU}{source}{}{sourcepointer-5}{}}
\newlabel{topology-theorem-tychonov@cref}{{[originaltheorem][0][]08ZU}{[1][0][]source}{}{}{}}
\makeatother
\title{The category of locally quasi-compact Hausdorff spaces has fibre products, a final object and quasi-compact fibre products of quasi-compact objects}
\author{The Stacks Project Authors}
\date{stacks-project:sites-cohomology.tex:local-snapshot-20261011; GNU FDL 1.2 or later}
\begin{document}
\maketitle
\noindent\textit{Editorial note.} All delivered mathematical text is the author's own, copied from the pinned Stacks Project chapter file sites-cohomology.tex; no statement or proof step was rewritten, reordered or removed.\par\medskip
\noindent\textit{Editorial note.} The counted proof is the author's complete argument: the proof is written out in full: the final object and Hausdorffness, the quasi-compact case through Tychonoff and closedness, and the general case through quasi-compact neighbourhoods.\par\medskip
\noindent\textit{Editorial note.} The label lemma-LC-basic names the counted statement, and the proof environment immediately following it in the source is delivered verbatim as its proof.\par\medskip
\noindent\textit{Editorial note.} The counted statement is carried by the author's own lemma environment in the chapter source.\par\medskip
\noindent\textit{Editorial note.} The statement and proof are copied from the pinned chapter file; the five cited inputs are kept as external pointers with their canonical Stacks tags.\par\medskip
\noindent\textit{Allowed standard prerequisites (editorial).}\begin{itemize}
\item Hausdorffness of products and subspaces. Topology, the section on Hausdorff spaces, giving the Hausdorff property of X x Y and of its subspaces. Supplies the Hausdorff property of the fibre product as a subspace of the product.
\item Tychonoff and closedness of the fibre product. Topology, Tychonoff's theorem and the lemma stating that the fibre product is closed in the product. Supplies quasi-compactness of the fibre product of quasi-compact objects.
\item Criterion for being locally quasi-compact. Topology, the lemma characterising locally quasi-compact Hausdorff spaces through quasi-compact neighbourhoods. Supplies the final step producing a quasi-compact neighbourhood of an arbitrary point.
\end{itemize}\medskip
\ifdefined\hypertarget\hypertarget{sourcepointer-1}{}\fi
\noindent\textit{Source reference.} Original lemma 005C: sources/stacks-chow/tags.txt, source lines 204--204. Prerequisite pointer invoked at proof line(s) 7494; it names a standard fact of a sibling chapter that the counted argument uses, and every step of the counted proof built on it is retained verbatim in the delivered body.\par\medskip
\ifdefined\hypertarget\hypertarget{sourcepointer-2}{}\fi
\noindent\textit{Source reference.} Original lemma 08ZH: sources/stacks-chow/tags.txt, source lines 11375--11375. Prerequisite pointer invoked at proof line(s) 7492; it names a standard fact of a sibling chapter that the counted argument uses, and every step of the counted proof built on it is retained verbatim in the delivered body.\par\medskip
\ifdefined\hypertarget\hypertarget{sourcepointer-3}{}\fi
\noindent\textit{Source reference.} Original lemma 08ZR: sources/stacks-chow/tags.txt, source lines 11384--11384. Prerequisite pointer invoked at proof line(s) 7502; it names a standard fact of a sibling chapter that the counted argument uses, and every step of the counted proof built on it is retained verbatim in the delivered body.\par\medskip
\ifdefined\hypertarget\hypertarget{sourcepointer-4}{}\fi
\noindent\textit{Source reference.} Original section 08ZD: sources/stacks-chow/tags.txt, source lines 11371--11371. Prerequisite pointer invoked at proof line(s) 7486; it names a standard fact of a sibling chapter that the counted argument uses, and every step of the counted proof built on it is retained verbatim in the delivered body.\par\medskip
\ifdefined\hypertarget\hypertarget{sourcepointer-5}{}\fi
\noindent\textit{Source reference.} Original theorem 08ZU: sources/stacks-chow/tags.txt, source lines 11387--11387. Prerequisite pointer invoked at proof line(s) 7490; it names a standard fact of a sibling chapter that the counted argument uses, and every step of the counted proof built on it is retained verbatim in the delivered body.\par\medskip
\begin{lemma}
\label{lemma-LC-basic}
The category $\textit{LC}$ has fibre products and a final object and hence
has arbitrary finite limits. Given morphisms $X \to Z$ and $Y \to Z$
in $\textit{LC}$ with
$X$ and $Y$ quasi-compact, then $X \times_Z Y$ is quasi-compact.
\end{lemma}

\begin{proof}
The final object is the singleton space. Given morphisms $X \to Z$ and
$Y \to Z$ of $\textit{LC}$ the fibre product $X \times_Z Y$ is
a subspace of $X \times Y$. Hence $X \times_Z Y$ is Hausdorff as
$X \times Y$ is Hausdorff by
Topology, Section \ref{topology-section-Hausdorff}.

\medskip\noindent
If $X$ and $Y$ are quasi-compact, then $X \times Y$ is quasi-compact by 
Topology, Theorem \ref{topology-theorem-tychonov}.
Since $X \times_Z Y$ is a closed subset of $X \times Y$
(Topology, Lemma \ref{topology-lemma-fibre-product-closed})
we find that $X \times_Z Y$ is quasi-compact by
Topology, Lemma \ref{topology-lemma-closed-in-quasi-compact}.

\medskip\noindent
Finally, returning to the general case, if $x \in X$ and $y \in Y$
we can pick quasi-compact neighbourhoods $x \in E \subset X$ and
$y \in F \subset Y$ and we find that $E \times_Z F$ is a quasi-compact
neighbourhood of $(x, y)$ by the result above. Thus $X \times_Z Y$
is an object of $\textit{LC}$ by
Topology, Lemma \ref{topology-lemma-locally-quasi-compact-Hausdorff}.
\end{proof}

\end{document}
